\subsection{多项式的构造}

令 A 为环 Z{[}X, Y{]}，它是具有整数系数、关于两族不定元 \(\mathbf{X} = (\mathbf{X}_{n})_{n \in \mathbb{N}}\) 和 \(\mathbf{Y} = (\mathbf{Y}_{n})_{n \in \mathbb{N}}\) 的多项式环。令 \(\theta\) 为 A 的自同态，由 \(\theta(\mathbf{X}_{n}) = \mathbf{X}_{n}^{p}\) 及 \(\theta(\mathbf{Y}_{n}) = \mathbf{Y}_{n}^{p}\)（对每个 \(n \in \mathbf{N}\)）定义。则 p 不是 A 中的零因子，且 A 中满足 \(\theta(a) \equiv a^{p} \mod p\) 的 a 的集合。A 是包含 X\_n 和 Y\_n 的 A 的子环，因而等于 A。

 根据第 2 号的命题 2 的 a) 和 c)，存在元素 \(\mathbf{S} = (\mathrm{S}_{n})_{n\in\mathbb{N}}\)、\(\mathbf{P} = (\mathrm{P}_{n})_{n\in\mathbb{N}}\)、\(\mathbf{I} = (\mathrm{I}_{n})_{n\in\mathbb{N}}\) 和 \(\mathbf{F} = (\mathrm{F}_{n})_{n\in\mathbb{N}}\)，它们属于 \(\mathrm{A}^{\mathbf{N}}\)，分别由下列等式刻画

\[
\left\{ \begin{array}{l} \Phi_ {\mathbf {A}} (\mathbf {S}) = \Phi_ {\mathbf {A}} (\mathbf {X}) + \Phi_ {\mathbf {A}} (\mathbf {Y}) \\ \Phi_ {\mathbf {A}} (\mathbf {P}) = \Phi_ {\mathbf {A}} (\mathbf {X}) \Phi_ {\mathbf {A}} (\mathbf {Y}) \\ \Phi_ {\mathbf {A}} (\mathbf {I}) = - \Phi_ {\mathbf {A}} (\mathbf {X}) \\ \Phi_ {\mathbf {A}} (\mathbf {F}) = f _ {\mathbf {A}} \big (\Phi_ {\mathbf {A}} (\mathbf {X}) \big). \end{array} \right.\tag{11}
\]

因此，A 中的元素 \(S_{n}\)、\(P_{n}\)、\(I_{n}\) 和 \(F_{n}\) 分别由下列公式刻画（其中 \(n\) 遍历 \(\mathbf{N}\)）：

\begin{enumerate}
\def\labelenumi{(\arabic{enumi})}
\setcounter{enumi}{11}
\tightlist
\item
\end{enumerate}

\[
\Phi_ {n} \left(\mathrm{S} _ {0}, \dots , \mathrm{S} _ {n}\right) = \Phi_ {n} \left(\mathrm{X} _ {0}, \dots , \mathrm{X} _ {n}\right) + \Phi_ {n} \left(\mathrm{Y} _ {0}, \dots , \mathrm{Y} _ {n}\right),\tag{13}
\]

\[
\Phi_ {n} \left(\mathrm{P} _ {0}, \dots , \mathrm{P} _ {n}\right) = \Phi_ {n} \left(\mathrm{X} _ {0}, \dots , \mathrm{X} _ {n}\right) \Phi_ {n} \left(\mathrm{Y} _ {0}, \dots , \mathrm{Y} _ {n}\right),\tag{14}
\]

\[
\Phi_ {n} (\mathrm{I} _ {0}, \dots , \mathrm{I} _ {n}) = - \Phi_ {n} (\mathrm{X} _ {0}, \dots , \mathrm{X} _ {n}),\tag{15}
\]

\[
\Phi_ {n} \left(\mathrm{F} _ {0}, \dots , \mathrm{F} _ {n}\right) = \Phi_ {n + 1} \left(\mathrm{X} _ {0}, \dots , \mathrm{X} _ {n + 1}\right).
\]

给 \(X_{n}\) 和 \(Y_{n}\) 赋予权 \(p^{n}\)，对每个 \(n \in \mathbf{N}\)。由第 2 号的注记推出下列断言：

\begin{enumerate}
\def\labelenumi{\alph{enumi})}
\item
  有 \(\mathbf{S}_n\in \mathbf{Z}[\mathbf{X}_0,\dots ,\mathbf{X}_n,\mathbf{Y}_0,\dots ,\mathbf{Y}_n]\)，且 \(\mathbf{S}_n\) 是权为 \(p^n\) 的等重多项式。
\item
  有 \(\mathbf{P}_n\in \mathbf{Z}[\mathrm{X}_0,\dots ,\mathrm{X}_n,\mathrm{Y}_0,\dots ,\mathrm{Y}_n]\)，且 \(\mathbf{P}_n\) 是权为 \(p^n\) 的等重多项式，分别关于族 \((\mathrm{X}_0,\dots ,\mathrm{X}_n)\) 和 \((\mathrm{Y}_0,\dots ,\mathrm{Y}_n)\)。
\item
  有 \(\mathbf{I}_n\in \mathbf{Z}[\mathrm{X}_0,\dots ,\mathrm{X}_n]\)，且 \(\mathbf{I}_n\) 是权为 \(p^n\) 的等重多项式。
\item
  有 \(\mathbf{F}_n\in \mathbf{Z}[\mathbf{X}_0,\dots ,\mathbf{X}_{n + 1}]\)，且 \(\mathbf{F}_n\) 是权为 \(p^{n + 1}\) 的等重多项式。
\end{enumerate}

公式 (2) 实际上允许逐步确定多项式 \(S_{n}\)、\(P_{n}\)、\(I_{n}\) 和 \(F_{n}\)。

例子。-- 1) 有

\[
\mathbf {S _ {0}} = \mathbf {X _ {0}} + \mathbf {Y _ {0}}
\]

\[
\mathrm{S} _ {1} = \mathrm{X} _ {1} + \mathrm{Y} _ {1} - \sum_ {i = 1} ^ {p - 1} \frac {1}{p} \binom {p} {i} \mathrm{X} _ {0} ^ {i} \mathrm{Y} _ {0} ^ {p - i}.
\]

此外，\(S_{n} - X_{n} - Y_{n}\) 属于环 \(Z[X_{0}, ..., X_{n-1}, Y_{0}, ..., Y_{n-1}]\)。2) 有

\[
\begin{array}{l} \mathbf {P} _ {0} = \mathbf {X} _ {0} \mathbf {Y} _ {0} \\ \mathbf {P} _ {1} = p \mathbf {X} _ {1} \mathbf {Y} _ {1} + \mathbf {X} _ {0} ^ {p} \mathbf {Y} _ {1} + \mathbf {X} _ {1} \mathbf {Y} _ {0} ^ {p}. \end{array}
\]

\begin{enumerate}
\def\labelenumi{\arabic{enumi})}
\setcounter{enumi}{2}
\tightlist
\item
  当 \(p \neq 2\) 时，有 \(I_n = -X_n\)。当 \(p = 2\) 时，有
\end{enumerate}

\[
\mathrm{I} _ {0} = - \mathrm{X} _ {0}
\]

\[
\mathrm{I} _ {1} = - \left(\mathrm{X} _ {0} ^ {2} + \mathrm{X} _ {1}\right)
\]

\[
\mathrm{I} _ {2} = - \mathrm{X} _ {0} ^ {4} - \mathrm{X} _ {0} ^ {2} \mathrm{X} _ {1} - \mathrm{X} _ {1} ^ {2} - \mathrm{X} _ {2}.
\]

\begin{enumerate}
\def\labelenumi{\arabic{enumi})}
\setcounter{enumi}{3}
\tightlist
\item
  有
\end{enumerate}

\[
\mathbf {F} _ {0} = \mathbf {X} _ {0} ^ {p} + p \mathbf {X} _ {1}
\]

\[
\mathrm{F} _ {1} = \mathrm{X} _ {1} ^ {p} + p \mathrm{X} _ {2} - \sum_ {i = 0} ^ {p - 1} \binom {p} {i} p ^ {p - i - 1} \mathrm{X} _ {0} ^ {p i} \mathrm{X} _ {1} ^ {p - i}.
\]

  由于 \(\Phi_n(F_0, \ldots, F_n) \equiv \Phi_n(X_0^p, \ldots, X_n^p) \pmod{p^{n+1}A}\) 对所有 \(n \in \mathbf{N}\) 成立（公式 (2) 和 (15)），根据命题 1(b) 可得 \(F_n \equiv X_n^p \pmod{pA}\) 对所有 \(n \in \mathbf{N}\) 成立。

 注记。--- 令 J 为满足 \(j \geqslant 1\) 的整数集。对于 J 的每个元素 \(j\)，定义多项式 \(\varphi_j\)，它属于 \(\mathbf{Z}[(\mathbf{X}_j)_{j \in \mathbf{J}}]\)，按下式给出

\[
\varphi_ {j} = \sum_ {d} d X _ {d} ^ {j / d},
\]

其中求和遍历 J 中整除 j 的元素。对于每个整数 \(n \geqslant 0\)，有

\[
\varphi_ {p ^ {n}} = \Phi_ {n} (\mathrm{X} _ {p ^ {0}},..., \mathrm{X} _ {p ^ {n}}).
\]

对于每个环 A 及每个元素 \(m\)（属于 \(\mathbf{J}\)），以 \(\varphi_{m}\) 表示从 \(A^{\mathbf{J}}\) 到 A 的映射，它将 \((a_{j})_{j\in\mathbf{J}}\) 送到 \(\varphi_{m}((a_{j})_{j\in\mathbf{J}})\)；以 \(\varphi_{A}\) 或简称 \(\varphi\) 表示从 \(A^{\mathbf{J}}\) 到自身的映射，它将 \(a = (a_{j})_{j\in\mathbf{J}}\) 送到 \((\varphi_{m}(a))_{m\in\mathbf{J}}\)。

令 \(\mathcal{A} = \mathbf{Z}[(X_j)_{j\in\mathbf{J}},(Y_j)_{j\in\mathbf{J}}]\) 为关于两族不定元 \(\mathbf{X} = (X_j)_{j\in\mathbf{J}}\) 和 \(\mathbf{Y} = (Y_j)_{j\in\mathbf{J}}\) 的整数系数多项式环。可以证明（p. 51, 习题 34），\(\mathcal{A}\) 中存在元素

\[
\boldsymbol {s} = (s _ {j}) _ {j \in \mathbf {J}}, \boldsymbol {p} = (p _ {j}) _ {j \in \mathbf {J}} \quad \text{且} \quad \boldsymbol {i} = (i _ {j}) _ {j \in \mathbf {J}},
\]

它们分别由下列等式刻画：

\[
\begin{array}{l} \varphi_ {\mathcal {A}} (s) = \varphi_ {\mathcal {A}} (\mathbf {X}) + \varphi_ {\mathcal {A}} (\mathbf {Y}) \\ \varphi_ {\mathcal {A}} (p) = \varphi_ {\mathcal {A}} (\mathbf {X}) \varphi_ {\mathcal {A}} (\mathbf {Y}) \\ \varphi_ {\mathcal {A}} (i) = - \varphi_ {\mathcal {A}} (\mathbf {X}). \end{array}
\]

